\begin{center}\textls[100]{d. Transzendente Zahlen. Asymptotische Gesetze und mittlere Werte.}\\\textls[100]{Bestimmung der Klassenanzahl.}\end{center}
\begin{center}\bfseries 26.\end{center}
\noindent\hspace*{\parindent}Um vollständig zu sein, muss noch erwähnt werden, dass \textsc{Gauss} auch der Frage nach der arithmetischen Beschaffenheit der Zahl \(\pi\) nahe getreten ist und einen eigenen Beweis für die Irrationalität der Tangenten rationaler Bögen gegeben hat (W. VIII, S. 27) nebst kritischen Bemerkungen über die voraufgegangenen bezüglichen Arbeiten von \textsc{Lambert} und \textsc{Legendre}\footnote{Die kritischen Bemerkungen, von denen im Text die Rede ist, werden von R. \textsc{Fricke} in seiner Bemerkung W. VIII, S. 29 erwähnt, sind aber nicht abgedruckt. Sie stehen auf mehreren Kleinoktavblättern und sind augenscheinlich früher geschrieben als der auf einem Grossquartblatte stehende, W. VIII, S. 27, 28 abgedruckte »Beweis« (beides in Ee, Kapsel 45). Von diesem »Beweis« findet sich übrigens auch noch eine von fremder Hand herrührende Abschrift, die am Schluss die Angabe enthält »\textsc{Gauss} 1850 August 20«. Die kritischen Bemerkungen beziehen sich auf die von \textsc{Legendre} in seinen \textit{Élémens de Géométrie} (S. 295 in der 7. Ausgabe von 1808) erwähnte Abhandlung von \textsc{Lambert} »in den Mem. de l’Ac. de Berlin 1761«. Gemeint ist offenbar die Abhandlung \textit{Mémoire sur quelques propriétés remarquables des quantités transcendantes circulaires et logarithmiques}, lu en 1767, Histoire de l’Académie, Berlin (1761) 1768, S. 265. \textsc{Gauss} sagt in bezug auf den dort von \textsc{Lambert} gegebenen Beweis für die Irrationalität von \(\pi\): »Dieser Beweis ist auf eine ziemlich schwerfällige und verworrene Art dargestellt, und man könnte sagen, dass die Denkschrift nicht sowohl einen klar entwickelten Beweis enthalte, als vielmehr nur das Material zu dem Beweise. Folgende Erläuterungen können dazu dienen, das Verständnis zu erleichtern«. In bezug auf diese Abhandlung \textsc{Lamberts} vergl. man den Aufsatz von A. \textsc{Pringsheim} »\textit{Über die ersten Beweise der Irrationalität von \(e\) und \(\pi\)}«, Münchener Sitzungsberichte \textbf{28}, 1898, S. 325, wo auch \textsc{Euler} als Vorgänger von \textsc{Lambert} zu seinem Rechte kommt.}. Desgleichen
\newpage\SourceLeaf{70}\ArithmeticHead{63}{R}{TRANSZENDENTE ZAHLEN.}
\noindent sind hier die umfangreichen Rechnungen und Tabellen (W. II, S. 477—495) zur Zyklotechnie zu nennen, die \textsc{Gauss} ausgeführt hat und welche den Zweck verfolgen, die Bögen, deren Kotangenten gegebene rationale Zahlen sind, genau zu berechnen. Auf Grund der Formel
\[
\operatorname{arc\,cotg}\frac{y}{a}=\operatorname{arc\,cotg}\frac{y+x}{a}+\operatorname{arc\,cotg}\frac{y+z}{a}
\]
für
\[
y^2+a^2=xz
\]
lassen sich die Bögen für kleine Kotangenten aus denen für grosse bestimmen, für welche die bekannten Reihen stärker konvergieren. So hat \textsc{Gauss} u. a. (s. W. II, S. 501) für \(\frac\pi4\) die beiden Ausdrücke
\[
\begin{aligned}
\frac\pi4&=12(18)+8(57)-5(239)\\
&=12(38)+20(57)+7(239)+24(268)
\end{aligned}
\]
gegeben, in denen \((18)=\operatorname{arc\,cotg}18\) usw. ist, und die wesentlich schneller zum Ziele führen als frühere, von \textsc{Machin}, \textsc{Euler}, \textsc{Vega} u. A. gegebene ähnliche Formeln.

Veröffentlicht hat \textsc{Gauss} weitere zahlentheoretische Arbeiten nicht, obwohl er zweifelsohne noch Manches gewusst hat, dessen Kenntnis für uns wertvoll gewesen sein würde. Wir haben nun aber, auf die Zeit der D. A. zurückgehend, in Anknüpfung an die artt. 301 und 302 derselben noch von einem andern Gebiete zu handeln, dessen Aufschliessung wir gewohnt sind, \textsc{Dirichlet}
\newpage\SourceLeaf{71}\ArithmeticHead{64}{L}{BACHMANN, ÜBER GAUSS’ ZAHLENTHEORETISCHE ARBEITEN.}
\noindent zuzusprechen, wenn wir von seinen »analytischen Methoden der Zahlentheorie« reden. Den selbständigen Leistungen dieses Forschers geschieht kein Abbruch, wenn man auch nach dem Zeugnis des handschriftlichen Nachlasses von \textsc{Gauss} anerkennen muss, dass vor \textsc{Dirichlet} schon \textsc{Gauss} auch in diesem Gebiete Pfadfinder gewesen ist.

Es handelt sich zunächst um asymptotische Gesetze der Zahlentheorie\footnote{Zum Folgenden vergl. W. X1, S. 11 ff., wo die im Text besprochenen Nachlassstücke abgedruckt sind; in den Bemerkungen S. 16 ff. findet man auch Hinweise auf die neuere Literatur.}. Hier hatte \textsc{Gauss}, wie erwähnt, mit der Abzählung der Primzahlen frühzeitig begonnen, sie nach 1811 von Zeit zu Zeit fortgesetzt und später durch \textsc{Goldschmidt} bis zum Umfange der Tafel W. II, S. 435 weiterführen lassen; auch hatte er bald schon den Integrallogarithmus als geeigneten asymptotischen Ausdruck für die Menge der Primzahlen unter einer gegebenen Grenze erkannt. In der Tat finden sich in \textsc{Gauss}’ Exemplar der \textit{logarithmischen Tafeln} von \textsc{Schulze} Aufzeichnungen vom Mai 1796 (vergl. Nr. 9 T. vom 31. Mai 1796), nach denen die Menge der Primzahlen unterhalb \(a\) asymptotisch durch \(\frac{a}{\log a}\), die Menge der Zahlen, welche aus zwei Primfaktoren zusammengesetzt sind, durch \(\log\log a\mathbin{.}\frac{a}{\log a}\) und »wahrscheinlich« die Menge derjenigen, welche aus drei Primfaktoren bestehen, durch \(\frac12(\log\log a)^2\mathbin{.}\frac{a}{\log a}\) usw. in infinitum ausgedrückt werde. Ferner gibt \textsc{Gauss} an derselben Stelle für die Menge der Zahlen unterhalb \(a\) ohne gleiche Faktoren den asymptotischen Ausdruck
\[
\frac{a}{\displaystyle\sum_n\frac1{n^2}},
\]
allgemeiner für die Menge derjenigen, welche höchstens \(\kappa=2,3,\ldots\) gleiche Faktoren haben, den Ausdruck
\[
\frac{a}{\displaystyle\sum_n\frac1{n^{\kappa+1}}}.
\]
Desgleichen findet er für die Summe der reziproken Primzahlen unterhalb \(x\) die Gleichung
\[
\sum\frac1p=\log\log x+V,
\]
\newpage\SourceLeaf{72}\ArithmeticHead{65}{R}{ASYMPTOTISCHE GESETZE.}\HistoricalFoot{X\textsubscript{2}.}{9}
\noindent wo \(V\) wahrscheinlich eine der Zahl \(1{,}266\) nahe Konstante sei und für das entsprechende Produkt
\[
\prod\frac1{1-\frac1p}
\]
den Ausdruck \(a\mathbin{.}\log x\), wo \(a\) eine der Zahl \(1{,}874\) nahe Konstante. Noch findet sich der Ausspruch: numerus factorum usque ad \(n\)
\[
(\log\log n+1)n.
\]

In \textsc{Gauss}’ Exemplar der \textsc{Lambert}schen \textit{Tabellen} wird für die Anzahl aller nur aus den Faktoren \(2,3,5,7\) bestehenden Zahlen unterhalb \(n\) der (nicht korrekte) asymptotische Ausdruck gegeben:
\[
\frac{\frac18(\log n)^4}{\log2\mathbin{.}\log3\mathbin{.}\log4\mathbin{.}\log5}.
\]
Desgleichen steht in einem im November 1799 begonnenen Hefte des Nachlasses (Ac) S. 19 der Satz, dass die Summe
\[
\sum_1^n\frac{\varphi(n)}{n^2},
\]
wo \(\varphi(n)\) die \textsc{Euler}sche Funktion, quam proxime durch \(\frac6{\pi^2}\log n+\mathrm{const.}\) ausgedrückt werde, während die Konstante etwa \(0{,}697413\) sei; und in einem andern im Juli 1800 begonnenen Hefte (Ae) wird S. 39 für die Anzahl \(\psi(M)\) der Teiler einer Zahl \(M\) die asymptotische Formel
\[
\sum\psi(M)=M(\log M+0{,}15443)+?
\]
gegeben. Die Notizen Nr. 14 und 31 des \textsc{Gauss}schen \textit{Tagebuchs} vom 20. Juni und 6. September 1796 sprechen ferner zwei asymptotische Gesetze aus, für die zuerst \textsc{Dirichlet} einen Beweis bekannt gemacht hat, und welche mit Benutzung des \textsc{Euler}schen Zeichens \(\int n\) für die Summe der Teiler einer Zahl \(n\) sowie der \textsc{Euler}schen Funktion \(\varphi(n)\) in den Formeln
\[
\frac{\int1+\int2+\ldots+\int n}{\frac{n(n+1)}2}=\frac{\pi^2}{6}
\]
d. i.
\[
\frac{\int1+\int2+\ldots+\int n}{n}=\frac{\pi^2}{12}\mathbin{.}n
\]
\newpage\SourceLeaf{73}\ArithmeticHead{66}{L}{BACHMANN, ÜBER GAUSS’ ZAHLENTHEORETISCHE ARBEITEN.}
\noindent beziehungsweise
\[
2\mathbin{.}\frac{\varphi(2)+\varphi(3)+\ldots+\varphi(n)}{n^2}=\frac6{\pi^2}
\]
ihren Ausdruck finden. Die letztere Formel hat \textsc{Gauss} auch auf einem aus den letzten Augusttagen des Jahres 1796 stammenden \textit{Exercitationes mathematicae} betitelten Blättchen (abgedruckt W. X1, S. 138, siehe besonders S. 140) aufgezeichnet.

Bei dem gänzlichen Fehlen jeder Andeutung einer für die Auffindung dieser Formeln dienenden Methode hat es fast den Anschein, als ob alle diese Ergebnisse nur induktiv gefunden seien, so schwer es andererseits zu denken ist, dass derartig verborgen liegende Beziehungen ohne den Anhalt bestimmter theoretischer Gesichtspunkte erkannt sein sollten.
